% ── Rozwiązanie: LCS -- instancja 2 ───────────────────────────────
\subsection{Najdłuższy wspólny podciąg (LCS) -- instancja 2}

Dane: $X = (G,\ A,\ T,\ T,\ A,\ C,\ A)$, $Y = (G,\ C,\ A,\ T,\ A,\ A,\ T)$, czyli
$m = 7$, $n = 7$. Tablicę kosztów $c$ wypełniamy wierszami regułą jak w~instancji~1,
a~równolegle w~tablicy $b$ zapisujemy kierunek ($\nwarrow$ -- dopasowanie,
$\uparrow$ -- góra, $\leftarrow$ -- lewo; remisy na korzyść góry). Komórki
wchodzące do wybranego LCS wyróżniono na~\textcolor{accentB}{pomarańczowo}.

\begin{aztable}
	\centering
	\renewcommand{\arraystretch}{1.2}
	\begin{minipage}{0.48\linewidth}
		\centering
		Tablica kosztów $c$: \\[0.5em]
		\begin{NiceTabular}{cc|c|ccccccc}[margin,hvlines]
			    & $j$   & 0     & 1   & 2   & 3   & 4   & 5   & 6   & 7   \\
			$i$ &       & $y_j$ & $G$ & $C$ & $A$ & $T$ & $A$ & $A$ & $T$ \\
			\hline
			0   & $x_i$ & 0     & 0   & 0   & 0   & 0   & 0   & 0   & 0   \\
			1   & $G$   & 0     & \color{accentB}1 & 1 & 1 & 1 & 1 & 1 & 1 \\
			2   & $A$   & 0     & 1   & 1   & \color{accentB}2 & 2 & 2 & 2 & 2 \\
			3   & $T$   & 0     & 1   & 1   & 2   & 3   & 3   & 3   & 3   \\
			4   & $T$   & 0     & 1   & 1   & 2   & \color{accentB}3 & 3 & 3 & 4 \\
			5   & $A$   & 0     & 1   & 1   & 2   & 3   & \color{accentB}4 & 4 & 4 \\
			6   & $C$   & 0     & 1   & 2   & 2   & 3   & 4   & 4   & 4   \\
			7   & $A$   & 0     & 1   & 2   & 3   & 3   & 4   & \color{accentB}5 & 5 \\
		\end{NiceTabular}
	\end{minipage}\hfill
	\begin{minipage}{0.48\linewidth}
		\centering
		Tablica strzałek $b$: \\[0.5em]
		\begin{NiceTabular}{cc|c|ccccccc}[margin,hvlines]
			    & $j$   & 0     & 1          & 2            & 3            & 4            & 5            & 6            & 7            \\
			$i$ &       & $y_j$ & $G$        & $C$          & $A$          & $T$          & $A$          & $A$          & $T$          \\
			\hline
			0   & $x_i$ &       &            &              &              &              &              &              &              \\
			1   & $G$   &       & \color{accentB}$\nwarrow$ & $\leftarrow$ & $\leftarrow$ & $\leftarrow$ & $\leftarrow$ & $\leftarrow$ & $\leftarrow$ \\
			2   & $A$   &       & $\uparrow$ & $\uparrow$   & \color{accentB}$\nwarrow$ & $\leftarrow$ & $\nwarrow$   & $\nwarrow$   & $\leftarrow$ \\
			3   & $T$   &       & $\uparrow$ & $\uparrow$   & $\uparrow$   & $\nwarrow$   & $\leftarrow$ & $\leftarrow$ & $\nwarrow$   \\
			4   & $T$   &       & $\uparrow$ & $\uparrow$   & $\uparrow$   & \color{accentB}$\nwarrow$ & $\uparrow$ & $\uparrow$ & $\nwarrow$ \\
			5   & $A$   &       & $\uparrow$ & $\uparrow$   & $\nwarrow$   & $\uparrow$   & \color{accentB}$\nwarrow$ & $\nwarrow$ & $\uparrow$ \\
			6   & $C$   &       & $\uparrow$ & $\nwarrow$   & $\uparrow$   & $\uparrow$   & $\uparrow$   & $\uparrow$   & $\uparrow$   \\
			7   & $A$   &       & $\uparrow$ & $\uparrow$   & $\nwarrow$   & $\uparrow$   & $\nwarrow$   & \color{accentB}$\nwarrow$ & $\leftarrow$ \\
		\end{NiceTabular}
	\end{minipage}
\end{aztable}

\begin{dygresja}
	Po jednej komórce na każdą gałąź reguły,
	$X = (x_1, \dots, x_7)$, $Y = (y_1, \dots, y_7)$:
	\begin{itemize}
		\item \textbf{Dopasowanie} ($\nwarrow$): $c(1, 1)$, $x_1 = G = y_1$, więc
		      $c(1, 1) = c(0, 0) + 1 = 1$.
		\item \textbf{Góra} ($\uparrow$): $c(2, 1)$, $x_2 = A \neq G = y_1$; sąsiedzi
		      $c(1, 1) = 1 \geq c(2, 0) = 0$, więc $c(2, 1) = 1$.
		\item \textbf{Lewo} ($\leftarrow$): $c(1, 2)$, $x_1 = G \neq C = y_2$; sąsiedzi
		      $c(0, 2) = 0 < c(1, 1) = 1$, więc $c(1, 2) = 1$.
		\item \textbf{Remis}: $c(2, 7)$, $A \neq T$, $c(1, 7) = c(2, 6) = 2$ -- warunek
		      $\geq$ spełniony, więc $\uparrow$.
	\end{itemize}
\end{dygresja}

\paragraph{Odtworzenie podciągu.}
Startujemy z~$(m, n) = (7, 7)$ i~idziemy za strzałkami; przy $\nwarrow$ schodzimy
po~przekątnej, dopisując $x_i$ do~LCS:
\[
	\begin{aligned}
		(7, 7) & \xrightarrow{\leftarrow} (7, 6) \xrightarrow{\nwarrow} (6, 5) \;[\,x_7 = A\,]
		\xrightarrow{\uparrow} (5, 5) \xrightarrow{\nwarrow} (4, 4) \;[\,x_5 = A\,] \\
		       & \xrightarrow{\nwarrow} (3, 3) \;[\,x_4 = T\,]
		\xrightarrow{\uparrow} (2, 3) \xrightarrow{\nwarrow} (1, 2) \;[\,x_2 = A\,]
		\xrightarrow{\leftarrow} (1, 1) \xrightarrow{\nwarrow} (0, 0) \;[\,x_1 = G\,].
	\end{aligned}
\]
Znaki dopisywane są \emph{po} powrocie z~rekurencji, więc w~kolejności od
najgłębszego: $G, A, T, A, A$. Stąd najdłuższy wspólny podciąg
\[
	\boxed{\text{LCS}(X, Y) = (G,\ A,\ T,\ A,\ A)}
\]
o~długości $c(7, 7) = 5$.
